\begin{table}[H]
\centering
\caption{Mean prediction-set sizes at 95\% nominal coverage in the
main simulation study.}
\label{tab:results-size-95}
\begingroup
\small
\setlength{\tabcolsep}{4pt}
\renewcommand{\arraystretch}{1.07}
\textbf{Panel A: Group 1 (all-state-1 histories)}\\[3pt]
\begin{tabular}{@{}rrrr@{}}
\toprule
$h$ & Targets & OCHS (all weights) & Randomized $i$-block \\
\midrule
3 & 706 & 3.262 & 6.788 \\
6 & 672 & 21.851 & 50.443 \\
9 & 639 & 118.851 & 392.574 \\
\bottomrule
\end{tabular}
\par\vspace{9pt}
\textbf{Panel B: Group 0 (other histories)}\\[3pt]
\begin{tabular}{@{}rrrrrr@{}}
\toprule
$h$ & Targets & OCHS $0.90$ & Count-only & OCHS $0.50$ & Randomized $i$-block \\
\midrule
3 & 4,294 & 1.707 & 1.690 & 2.333 & 6.535 \\
6 & 4,328 & 9.643 & 9.834 & 17.546 & 47.587 \\
9 & 4,361 & 37.852 & 36.217 & 114.871 & 355.473 \\
\bottomrule
\end{tabular}
\par\vspace{5pt}
\begin{minipage}{0.97\linewidth}
\footnotesize
Targets is the number of replicates in the indicated history group
out of 5,000 at each horizon. All three OCHS specifications use the
same augmented transition-count score in group~1. In Panel B,
OCHS headings indicate $\lambda$; count-only uses $\lambda=1$.
The $i$-block benchmark uses randomized ties. Mean sizes include
empty prediction sets.
\end{minipage}
\endgroup
\end{table}
